% Canonical Paper II source.
This appendix records the coordinate calculations behind
Section~\ref{sec:contact-analysis}.  They concern only the three-dimensional
contact state space $\mathcal C$; none of the fields below is a surface frame
field $X_u$ or $X_v$.

\subsection{Contact volume, brackets, and divergence}

From
\[
 \alpha=da-\mu\,du-\lambda a\,dv,
 \qquad d\alpha=-\lambda\,da\wedge dv,
\]
one obtains
\begin{equation}\label{eq:app-contact-volume-signed}
 \alpha\wedge d\alpha
 =-\mu\lambda\,du\wedge dv\wedge da.
\end{equation}
Thus $\alpha$ is a contact form precisely when $\mu\lambda\ne0$, and its
positive contact density is
$d\nu_{\mathrm c}=|\mu\lambda|\,du\,dv\,da$.

For completeness, the bracket calculation is
\begin{align}
 [D_u,D_v]
 &=[\partial_u+\mu\partial_a,
       \partial_v+\lambda a\partial_a]\notag\\
 &=\mu\,\partial_a(\lambda a)\partial_a
 =\mu\lambda T.\label{eq:app-contact-bracket}
\end{align}
In particular, $D_u,D_v,[D_u,D_v]$ have coefficient matrix
\[
 \begin{pmatrix}
  1&0&0\\
  0&1&0\\
  \mu&\lambda a&\mu\lambda
 \end{pmatrix}
\]
relative to $\partial_u,\partial_v,\partial_a$, and its determinant is
$\mu\lambda$.  This also verifies the step-two bracket condition directly.

Because $d\nu_{\mathrm c}$ has constant coordinate density,
\begin{equation}\label{eq:app-contact-coordinate-divergence}
 \diver_{\nu_{\mathrm c}}D_u
 =\partial_u(1)+\partial_a(\mu)=0,
 \qquad
 \diver_{\nu_{\mathrm c}}D_v
 =\partial_v(1)+\partial_a(\lambda a)=\lambda.
\end{equation}
For $f,g\in C_c^\infty(\mathcal C)$ and a real vector field $Y$,
integration of $\diver(f\overline gY)$ gives
\begin{equation}\label{eq:app-vector-field-adjoint}
 \int(Yf)\overline g\,d\nu_{\mathrm c}
 =\int f\,\overline{(-Y-\diver_{\nu_{\mathrm c}}Y)g}
      \,d\nu_{\mathrm c}.
\end{equation}
Equations~\eqref{eq:app-contact-coordinate-divergence} and
\eqref{eq:app-vector-field-adjoint} prove the adjoint formulas in
Proposition~\ref{prop:contact-divergence-adjoints}.

\subsection{Expansion of the variational operator}

The two squares expand as
\begin{align*}
 D_u^2
 &=\partial_u^2+2\mu\partial_u\partial_a+\mu^2\partial_a^2,\\
 D_v^2
 &=\partial_v^2+2\lambda a\partial_v\partial_a
   +\lambda^2a^2\partial_a^2+\lambda^2a\partial_a.
\end{align*}
Consequently the raw sum of squares is
\begin{equation}\label{eq:app-contact-raw-expanded}
 Q_{\mathcal C}
 =\partial_u^2+\partial_v^2
  +2\mu\partial_u\partial_a+2\lambda a\partial_v\partial_a
  +(\mu^2+\lambda^2a^2)\partial_a^2
  +\lambda^2a\partial_a,
\end{equation}
whereas the divergence-of-gradient expression is
\begin{equation}\label{eq:app-contact-variational-expanded}
 \Delta_{\mathcal C}
 =Q_{\mathcal C}+\lambda\partial_v+\lambda^2a\partial_a.
\end{equation}
The last two terms are the contact-volume drift.  In particular,
$Q_{\mathcal C}$ and $\Delta_{\mathcal C}$ are distinct when $\lambda\ne0$.

For compactly supported $f$,
\begin{align}
 \langle-\Delta_{\mathcal C}f,f\rangle_{\nu_{\mathrm c}}
 &=\langle D_u^*D_uf+D_v^*D_vf,f\rangle_{\nu_{\mathrm c}}\notag\\
 &=\|D_uf\|_{L^2(\nu_{\mathrm c})}^2
   +\|D_vf\|_{L^2(\nu_{\mathrm c})}^2.
 \label{eq:app-contact-energy}
\end{align}
This calculation identifies the differential expression represented by the
closed energy form; it does not by itself assert essential self-adjointness of
the minimal differential operator.

\subsection{Raw and adjoint CR products}

Writing $\dbarC=\frac12(D_u+iD_v)$ and
$\dC=\frac12(D_u-iD_v)$, direct multiplication gives
\begin{align}
 4\dC\dbarC
 &=(D_u-iD_v)(D_u+iD_v)
   =Q_{\mathcal C}+i[D_u,D_v]
   =Q_{\mathcal C}+i\mu\lambda T,
 \label{eq:app-contact-raw-product-plus}\\
 4\dbarC\dC
 &=(D_u+iD_v)(D_u-iD_v)
   =Q_{\mathcal C}-i[D_u,D_v]
   =Q_{\mathcal C}-i\mu\lambda T.
 \label{eq:app-contact-raw-product-minus}
\end{align}
These are algebraic identities and do not depend on a measure.

The contact-volume adjoints do depend on the measure.  Conjugating scalar
coefficients and using $D_u^*=-D_u$ and $D_v^*=-D_v-\lambda$ yields
\[
 \dbarC^*=\frac12(-D_u+iD_v+i\lambda),
 \qquad
 \dC^*=\frac12(-D_u-iD_v-i\lambda).
\]
The first energy product is
\begin{align}
 4\dbarC^*\dbarC
 &=(-D_u+iD_v+i\lambda)(D_u+iD_v)\notag\\
 &=-D_u^2-D_v^2-\lambda D_v
   +i\lambda D_u-i[D_u,D_v]\notag\\
 &=A_{\mathcal C}+i\lambda D_u-i\mu\lambda T.
 \label{eq:app-contact-adjoint-product-dbar}
\end{align}
Similarly,
\begin{align}
 4\dC^*\dC
 &=(-D_u-iD_v-i\lambda)(D_u-iD_v)\notag\\
 &=-D_u^2-D_v^2-\lambda D_v
   -i\lambda D_u+i[D_u,D_v]\notag\\
 &=A_{\mathcal C}-i\lambda D_u+i\mu\lambda T.
 \label{eq:app-contact-adjoint-product-d}
\end{align}
Since $\Reeb=-\mu^{-1}\partial_u$ and
$D_u=\partial_u+\mu T$, these formulas are equivalently
\begin{equation}\label{eq:app-contact-adjoint-products-Reeb}
 4\dbarC^*\dbarC=A_{\mathcal C}-i\mu\lambda\Reeb,
 \qquad
 4\dC^*\dC=A_{\mathcal C}+i\mu\lambda\Reeb.
\end{equation}
The difference between
Equations~\eqref{eq:app-contact-raw-product-plus}--\eqref{eq:app-contact-raw-product-minus}
and
Equations~\eqref{eq:app-contact-adjoint-product-dbar}--\eqref{eq:app-contact-adjoint-product-d}
is therefore not a sign convention: the latter products incorporate the
contact-volume adjoints.

\subsection{The balanced-measure conjugation}

Let $d\nu_{\mathrm b}=e^{-\lambda v}d\nu_{\mathrm c}$.  For a weighted
density $e^\phi d\nu_{\mathrm c}$,
\begin{equation}\label{eq:app-weighted-divergence}
 \diver_{e^\phi\nu_{\mathrm c}}Y
 =\diver_{\nu_{\mathrm c}}Y+Y\phi.
\end{equation}
Taking $\phi=-\lambda v$ gives
$D_u\phi=0$ and $D_v\phi=-\lambda$, so both horizontal divergences vanish
for $d\nu_{\mathrm b}$.  Hence the balanced-measure energy operator is the
formal expression $-Q_{\mathcal C}$.

The unitary map
\[
 U:L^2(d\nu_{\mathrm b})\longrightarrow L^2(d\nu_{\mathrm c}),
 \qquad Uf=e^{-\lambda v/2}f,
\]
satisfies, on test functions,
\begin{align}
 UD_uU^{-1}f&=D_uf,\label{eq:app-gauge-Du}\\
 UD_vU^{-1}f&=D_vf+\frac{\lambda}{2}f.
 \label{eq:app-gauge-Dv}
\end{align}
It follows that
\begin{align}
 U(-Q_{\mathcal C})U^{-1}
 &=-D_u^2-\left(D_v+\frac{\lambda}{2}\right)^2\notag\\
 &=-D_u^2-D_v^2-\lambda D_v-\frac{\lambda^2}{4}\notag\\
 &=A_{\mathcal C}-\frac{\lambda^2}{4}.
 \label{eq:app-gauge-energy}
\end{align}
Thus changing the measure and applying its unitary gauge produces the constant
zeroth-order shift in Equation~\eqref{eq:gauge-sum-squares}; it does not identify the
two unitarily conjugate differential expressions without that shift.

\subsection{Logarithmic first integrals}

Put $W(a)=\mu+i\lambda a$.  Since $\mu\ne0$, the image of $W$ is contained
in the half-plane $\Omega_\mu$ used in
Proposition~\ref{prop:contact-logarithmic-first-integrals}; therefore the chosen
branch $\mathrm{Log}_\mu W$ is smooth for every real $a$.  For
\[
 z=u+iv,
 \qquad
 \zeta=u+\frac{i}{\lambda}\mathrm{Log}_\mu W(a),
\]
one has
\begin{align}
 D_uz&=1,&D_vz&=i,\label{eq:app-z-derivatives}\\
 \partial_a\zeta&=-\frac1W,&
 D_u\zeta&=1-\frac{\mu}{W}=\frac{i\lambda a}{W},&
 D_v\zeta&=-\frac{\lambda a}{W}.
 \label{eq:app-zeta-derivatives}
\end{align}
Hence $D_uz+iD_vz=0$ and
$D_u\zeta+iD_v\zeta=0$.  This proves the CR assertions without treating
$\zeta$ as a real coordinate on $\mathcal C$.  Indeed, the useful global
object is the rank-three map $(z,\zeta):\mathcal C\to\C^2$.

For $\rho,\sigma\in\C$, the same branch convention gives
\begin{align}
 e^{\rho z+\sigma\zeta}
 &=e^{(\rho+\sigma)u+i\rho v}
   \exp\left(\frac{i\sigma}{\lambda}\mathrm{Log}_\mu W(a)\right)\notag\\
 &=e^{(\rho+\sigma)u+i\rho v}
   W(a)^{i\sigma/\lambda}.
 \label{eq:app-contact-mode-expanded}
\end{align}
This is a pointwise CR family.  The calculation supplies no Hilbert-space
orthogonality or completeness statement.

\subsection{Filtered affine--Appell calculations}

Let $r=e^{\lambda v}$ and $B_n=r^{-n}a^n$.  Since
$D_vr=\lambda r$, $D_va=\lambda a$, and $D_ua=\mu$, one finds
\begin{align}
 D_uB_n&=\mu n r^{-1}B_{n-1},\label{eq:app-Appell-Du}\\
 D_vB_n&=0,\label{eq:app-Appell-Dv}\\
 TB_n&=nr^{-1}B_{n-1}.
 \label{eq:app-Appell-T}
\end{align}
If $P\in C^\infty(\R^2_{(u,v)},\C)$, Leibniz' rule therefore gives the
module actions in Proposition~\ref{prop:contact-filtered-appell-module}.
The factor $r^{-1}$ is a smooth coefficient, but it is not a complex scalar;
this is exactly why the correct structure is a filtered
$C^\infty(\R^2_{(u,v)},\C)$-module rather than a finite-dimensional basis.

There is nevertheless a finite algebraic antiderivative for coefficients of
finite $u$-degree.  Write
$(n+1)_0=1$ and
$(n+1)_k=(n+1)(n+2)\cdots(n+k)$ for $k\ge1$.

\begin{proposition}[Finite upward sweep]\label{prop:app-contact-upward-sweep}
Let $P\in C^\infty(\R^2_{(u,v)},\C)$ satisfy
$\partial_u^{m+1}P=0$ for some $m\ge0$.  For $n\ge0$, define
\begin{equation}\label{eq:app-contact-upward-sweep}
 \mathcal I_u(PB_n)
 =\sum_{k=0}^{m}(-1)^k
  \frac{r^{k+1}\partial_u^kP}
       {\mu^{k+1}(n+1)_{k+1}}B_{n+k+1}.
\end{equation}
Then $\mathcal I_u(PB_n)$ belongs to $\mathcal M_{n+m+1}$ and
\begin{equation}\label{eq:app-contact-upward-sweep-identity}
 D_u\mathcal I_u(PB_n)=PB_n.
\end{equation}
\end{proposition}

\begin{proof}
Denote the $k$th summand in
Equation~\eqref{eq:app-contact-upward-sweep} by $S_k$.  Because $D_ur=0$,
Leibniz' rule and Equation~\eqref{eq:app-Appell-Du} give
\begin{align*}
 D_uS_k
 &=(-1)^k
   \frac{r^{k+1}\partial_u^{k+1}P}
        {\mu^{k+1}(n+1)_{k+1}}B_{n+k+1}\\
 &\quad+(-1)^k
   \frac{r^k\partial_u^kP}
        {\mu^k(n+1)_k}B_{n+k}.
\end{align*}
The second term for $k=0$ is $PB_n$.  For
$0\le k<m$, the first term for $k$ cancels the second term for $k+1$.
The only remaining first term is the one with $k=m$, and it vanishes because
$\partial_u^{m+1}P=0$.  This proves
Equation~\eqref{eq:app-contact-upward-sweep-identity}.
\end{proof}

The operator $\mathcal I_u$ in this proposition is only a finite algebraic
right inverse on the stated polynomial-in-$u$ coefficient class.  It is not a
bounded inverse on either contact $L^2$ space and gives no density, basis, or
completeness theorem for the filtered module.
